\block{\strut 2. Approach} {

    \vspace{0.4cm}

    \textbf{\underline{Workflow overview}}

    \vspace{0.8cm}

    \begin{center}
        \begin{tikzpicture}[x=\linewidth,y=8.2cm]
            \fill[approachOne] (0.00,0.06) -- (0.24,0.06) -- (0.32,0.50) -- (0.24,0.94) -- (0.00,0.94) -- cycle;
            \fill[approachTwo] (0.29,0.06) -- (0.55,0.06) -- (0.63,0.50) -- (0.55,0.94) -- (0.29,0.94) -- (0.37,0.50) -- cycle;
            \fill[approachThree] (0.60,0.06) -- (0.88,0.06) -- (0.96,0.50) -- (0.88,0.94) -- (0.60,0.94) -- (0.68,0.50) -- cycle;

            \node[align=left, text width=0.17\linewidth, font=\bfseries] at (0.13,0.52) {Generate\\nanoalloy\\configurations};
            \node[align=left, text width=0.18\linewidth, font=\bfseries] at (0.46,0.52) {Simulate HRTEM\\images and\\extract structural\\parameters};
            \node[align=left, text width=0.18\linewidth, font=\bfseries] at (0.79,0.52) {Train and apply\\deep-learning\\models};
        \end{tikzpicture}
    \end{center}

    \vspace{0.4cm}

    \textbf{MD simulations + multislice image generation + deep learning} provide a direct route from \textbf{synthetic data} to the interpretation of \textbf{experimental HRTEM images}.

}
